O engajamento dos estudantes é um componente crucial em ambientes de aprendizagem e um preditor significativo para a retenção e o sucesso acadêmico \cite{yilmazEngagementDistanceEducation2020}. Este construto possui uma forte relação com a aprendizagem e a satisfação discente, consolidando-se como uma tendência em pesquisas educacionais mundiais sob diversas perspectivas teóricas e metodológicas \cite{ortegaMensuracaoEngajamentoOnline2022}.

De forma ampla, o engajamento é compreendido como o tempo e a energia que os estudantes dedicam a atividades educacionalmente sólidas dentro e fora da sala de aula, e as políticas e práticas que as instituições usam para induzir os estudantes a participar dessas atividades \cite{kuhWhatWereLearning2003}. Envolve também um esforço diligente e uma participação ativa no processo de aprendizagem (Christenson, 2012 \textit{apud} \citeauthor{henrieMeasuringStudentEngagement2015}, \citeyear{henrieMeasuringStudentEngagement2015}). Engajar os estudantes é uma tarefa desafiadora, tanto na educação presencial quanto a distância \cite{yilmazEngagementDistanceEducation2020}.

Com a expansão da aprendizagem \textit{online}, surge um novo âmbito de interesse para a análise do engajamento: o \textit{engajamento online}. No contexto do aprendizado \textit{online}, o engajamento do aluno é um indicador relevante para avaliar sua aprendizagem, sendo seu nível um prognóstico do êxito educacional \cite{ortegaMensuracaoEngajamentoOnline2022}. Nessa modalidade, os estudantes frequentemente precisam ser autodirigidos e engajados em sua própria aprendizagem, devido à menor quantidade de interação com instrutores e colegas \cite{langePascarellaTerenzin20052014}. Quanto mais engajado está o estudante, maior a possibilidade de direcionar sua energia para a aprendizagem, potencialmente levando a resultados que ampliam ainda mais seus níveis de engajamento \cite{ortegaMensuracaoEngajamentoOnline2022}.

O número de estudos sobre engajamento em ambientes de educação a distância tem aumentado dramaticamente. Não é surpreendente que essa tendência de pesquisa esteja em ascensão, uma vez que, como apontam \citeonline{cochranRoleStudentCharacteristics2014}, as taxas de evasão são tradicionalmente mais altas na educação a distância do que no ensino tradicional e o engajamento é um fator determinante para a retenção discente \cite{yilmazEngagementDistanceEducation2020}. Nesse contexto, os recursos e ambientes de aprendizagem \textit{online} podem ser ajustados para atender às necessidades de distintos grupos de estudantes, personalizando a experiência educacional \cite{ortegaMensuracaoEngajamentoOnline2022}.

Para determinar o impacto de práticas instrucionais inovadoras na aprendizagem, são necessárias medidas úteis e válidas para mensurar o engajamento \cite{henrieMeasuringStudentEngagement2015}. A literatura na área tende a focar em medidas de autorrelato, particularmente escalas quantitativas. Sintetizando, a maioria dos artigos possui abordagem quantitativa e utiliza questionários para coleta de dados, em sua maioria com perguntas fechadas em escalas do tipo Likert de 5 pontos \cite{ortegaMensuracaoEngajamentoOnline2022}. No entanto, sugere-se que as bases teóricas que justificam a definição de cada subconstruto sejam estipuladas com clareza, especialmente considerando as características próprias dos ambientes \textit{online} \cite{ortegaMensuracaoEngajamentoOnline2022}.

Além dos métodos tradicionais de autorrelato, a tecnologia oferece novas formas de medir o engajamento de maneira escalável e minimamente disruptiva para a aprendizagem. Isso inclui a utilização de dados gerados por computador sobre a atividade do usuário em um sistema de aprendizagem constitui uma das novas possibilidades abertas pela tecnologia \cite{dmelloAutoTutorAffectiveAutotutor2013}. Essa abordagem amplia as frentes metodológicas para a pesquisa na área. Contudo, a busca por instrumentos que avaliem especificamente o \textit{engajamento online} nem sempre se revela plenamente no escopo de todos os questionários analisados, indicando uma área para desenvolvimento futuro \cite{ortegaMensuracaoEngajamentoOnline2022}.